\BourbakiExerciseGroup

\begin{enumerate}
\def\labelenumi{\arabic{enumi})}
\tightlist
\item
称点 \(a \in \mathbb{R}\) 为子集 \(\mathbf{A}\) （ \(\mathbb{R}\) 的）的左附着点，如果它属于 \(\mathbf{A}\) 的闭包，并且存在一个区间 \(]a, b[ (a < b)\) ，它不含 \(\mathbf{A}\) 的任何点。试证：子集 \(\mathbf{A}\) （ \(\mathbb{R}\) 的）的左附着点所成的集是可数的（在左附着点所成的集与一个由两两不相交的开区间组成的集之间建立一一对应）。由此试证 \(\mathbb{R}\) 的每个良序子集都是可数的。
\end{enumerate}

¶ 2) 设 \(\mathbf{A}\) 为 \(\mathbb{R}\) 中一个可数稠密子集，则它不是闭集。{[}若 \((a_{n})\) 是把 \(\mathbf{A}\) 的点按某个次序排列所得的序列，试定义一个区间序列 \([b_{n}, c_{n}]\) ，使得 \(b_{n-1} < b_{n} < c_{n} < c_{n-1}\) 对一切 \(n\) 成立，并且使得 \([b_{n}, c_{n}]\) 不含任何点 \(a_{k}\) ，其指标满足 \(k \leqslant n\) ；然后利用定理 2.{]}由此推出 \(\mathbb{R}\) 是不可数的（参看 § 8，第 6 节）。

\begin{enumerate}
\def\labelenumi{\arabic{enumi})}
\setcounter{enumi}{2}
\item
设 \((\mathbf{I}_n)\) 是 \(\mathbf{R}\) 中非空开区间的一个无穷序列，满足 \(\overline{\mathbf{I}}_n \cap \overline{\mathbf{I}}_m = \emptyset\) （ \(m \neq n\) ）。试证：诸 \(\mathbf{I}_n\) 的并的余集是一个完备集（第一章，§ 2，第 6 节）；特别地，Cantor 三分集是完备集。
\item
与 Cantor 三分集邻接的各区间的长度之和等于 1。试定义 {[}0, 1{]} 的一个全不连通完备子集，使得与 A 邻接且含于 {[}0, 1{]} 内的各区间的长度之和等于给定的数 m，它满足 \(0 < m \leqslant 1\) 。
\item
设 \(l\) 是一个数 \(> 0\) ，并假设对每个点 \(x \in \mathbf{R}\) 对应着一个开区间 \(\mathbf{I}(x)\) ，它以 \(x\) 为中心，长度 \(\leqslant l\) 。试证：每个紧区间 \([a, b]\) 都可被有限多个区间 \(\mathbf{I}(x_i)\) 覆盖，这些区间的长度之和 \(\leqslant l + 2(b - a)\) 。（试证：若命题对每个区间 \([a, x]\) （满足 \(a \leqslant x < c\) ）为真，则存在 \(d > c\) ，使它对每个区间 \([a, y]\) （满足 \(a \leqslant y < d\) ）为真。）若 \(l = (b - a)/n\) ，其中 \(n\) 是一个整数 \(\geqslant 1\) ，试证该结果不能再改进。
\end{enumerate}

¶ 6) a) 设 X 是一个非空全序集，赋予拓扑 \(\mathcal{C}_{0}(X)\) （第一章，§ 2，习题 5）。试证：X 是紧的当且仅当 X 的每个子集都有上确界，换言之，当且仅当 X 是一个完备格（集合论，第三章，§ 1，习题 11）。（为证条件必要，可按定理 3 的方式论证；为证充分性，考虑 X 上的一个滤子 \(\mathfrak{F}\) ，并证明：若 A 是 \(\mathfrak{F}\) 中诸集的下确界所成的集，则 A 的上确界是 \(\mathfrak{F}\) 的一个聚点。）

\begin{enumerate}
\def\labelenumi{\alph{enumi})}
\setcounter{enumi}{1}
\tightlist
\item
试给出一个完备格 X 的例子，它在拓扑 \(\mathcal{C}_{0}(\mathbf{X})\) 下不是紧的。
\end{enumerate}

¶ 7) 设 X 是一个非空全序集，赋予拓扑 \(\mathfrak{G}_{\mathfrak{0}}(\mathbf{X})\) 。

\begin{enumerate}
\def\labelenumi{\alph{enumi})}
\tightlist
\item
若 X 是连通的，试证它具有下列两个性质：
\end{enumerate}

α) X 的每个非空且上有界的子集都有上确界（若 A 非空且上有界，B 是 A 的上界所成的集，C 是 B 的下界所成的集，试证 \(B \cup C = X\) ，并且 B 与 C 都是闭集）。

β) 集 X 没有间隙，换言之（集合论，第三章，§ 1，习题 18），X 的每个开区间 {]}a, b{[} （a \textless{} b）都非空（利用定理 4 的论证）。

\begin{enumerate}
\def\labelenumi{\alph{enumi})}
\setcounter{enumi}{1}
\item
反之，试证若 X 具有性质 \(\alpha\) 与 \(\beta\) ），则 X 是连通的。（先证明每个闭区间 \([a, b]\) 都是连通的：设该区间是两个非空闭集 A 与 B 的不交并，且 \(a \in A\) ，考察 B 的下确界，由此得出矛盾。）
\item
试证：若 X 连通，则它是局部紧且局部连通的，并且 X 的仅有的连通子集就是区间（有界或无界）。
\end{enumerate}

\begin{enumerate}
\def\labelenumi{\arabic{enumi})}
\setcounter{enumi}{7}
\tightlist
\item
设 X 与 Y 是两个全序集，分别赋予拓扑 \(\mathcal{G}_{0}(X)\) 与 \(\mathcal{G}_{0}(Y)\) 。
\end{enumerate}

\begin{enumerate}
\def\labelenumi{\alph{enumi})}
\tightlist
\item
设 \(\mathbf{X}\) 连通，且设 \(f: \mathbf{X} \to \mathbf{Y}\) 是一个连续映射。试证：设 \(x, y\) 为 \(\mathbf{X}\) 中满足 \(x < y\) 的任意两点，则每个 \(z \in \mathbf{Y}\) ，只要它属于以 \(f(x)\) 与 \(f(y)\) 为端点的闭区间，就属于 \(f(\mathbf{X})\)
\end{enumerate}

（利用习题 7）。由此试证：\(f\) 是 \(\mathbf{X}\) 到 \(f(\mathbf{X})\) 上的同胚，当且仅当 \(f\) 连续且严格单调。

\begin{enumerate}
\def\labelenumi{\alph{enumi})}
\setcounter{enumi}{1}
\tightlist
\item
试给出有理直线 Q 到其自身的一个同胚的例子，它不是严格单调的。
\end{enumerate}

\begin{enumerate}
\def\labelenumi{\arabic{enumi})}
\setcounter{enumi}{8}
\tightlist
\item
  \begin{enumerate}
  \def\labelenumii{\alph{enumii})}
  \tightlist
  \item
设 X 是一个无间隙的可数全序集（习题 7）。试证：存在一个严格递增的双射映射 \(\varphi\) ，它把 X 映到以 o 与 i 为端点的 Q 的四个区间之一上。{[}把 X 的元素与 Q 的适当区间的点排成序列 \((a_n), (b_n)\) ，并用归纳法定义 \(\varphi\) 。{]}若 X 赋予拓扑 \(\tilde{C}_0(X)\) ，则 \(\varphi\) 是 X 到 \(\varphi(X)\) 上的同胚。
  \end{enumerate}
\end{enumerate}

\begin{enumerate}
\def\labelenumi{\alph{enumi})}
\setcounter{enumi}{1}
\item
试由 a) 推出：R 的开区间 {]}a, b{[} 的每个可数稠密子集都同胚于 Q。
\item
试由 a) 推出：若 X 是任一赋予拓扑 \(\mathfrak{G}_{0}(\mathbf{X})\) 的可数全序集，则存在一个严格递增的同胚，把 X 映到 Q 的一个子空间上{[}把 X 嵌入一个无间隙的可数全序空间 \(X'\) ，使得 \(\mathfrak{G}_{0}(\mathbf{X})\) 由 \(\mathfrak{G}_{0}(\mathbf{X}')\) 诱导{]}。
\end{enumerate}

\begin{enumerate}
\def\labelenumi{\arabic{enumi})}
\setcounter{enumi}{9}
\tightlist
\item
试证：Cantor 三分集 K 的每个可数子集，只要它在 K 中稠密且不含与 K 邻接的区间的任何端点，就同胚于有理直线 Q（利用习题 9）。
\end{enumerate}

¶ 11) a) 设 X 是带拓扑 \(\mathfrak{G}_{0}(\mathbf{X})\) 的全序集。若 X 连通且有一个可数稠密子集 A，试证：存在一个（由习题 8 知为严格单调的）同胚，把 X 映到以 o, i 为端点的 R 的区间之一上，并且把 A 映到 Q 与该区间的交上（利用习题 9 与 7）。

\begin{enumerate}
\def\labelenumi{\alph{enumi})}
\setcounter{enumi}{1}
\item
试由 a) 推出：若 B 是 R 的一个子集，它的余集在 R 中稠密，则存在 R 到其自身的同胚，它把 B 映到无理数集 CQ 的一个子集上（考虑一个含于 CB 中的 R 的可数稠密子集 A）。
\item
在按字典序排成全序的集合 \(\mathbf{R} \times \{0, 1\}\) 中，用 \(\mathbf{R}'\) 表示 \(\mathbf{Q} \times \{1\}\) 的余集。试证：\(\mathbf{R}'\) 赋予拓扑 \(\mathfrak{G}_0(\mathbf{R}')\) 后是局部紧且全不连通的，并且含有一个可数稠密子集 \(\mathbf{Q}'\) ，它同胚于 \(\mathbf{Q}\) ，但在 \(\mathbf{R}'\) 的非空开区间上诱导的拓扑不具有可数基；特别地，这样的区间不能同胚于 \(\mathbf{R}\) 的任何子集。
\end{enumerate}

¶ 12) a) 设 A 是一个良序集，用 I 表示 R 的区间 {[}0, 1{[} 。试证：按字典序排成全序并赋予拓扑 \(\mathfrak{G}_{0}(X)\) 的集合 X = A × I 是连通的（习题 7）；因此存在在拓扑 \(\mathfrak{G}_{0}(\mathbf{X})\) 下连通且具有任意势的全序集 X{[}参看 § 4，习题 7 b){]}。

\begin{enumerate}
\def\labelenumi{\alph{enumi})}
\setcounter{enumi}{1}
\tightlist
\item
取 A 为一个不可数的良序集，其中每个线段 \(\leftarrow, t]\) 都是可数的：相应的拓扑空间 X 称为 Alexandroff 半直线。试证：对 X 的每个区间 \([a, b]\) （a \textless{} b），存在一个严格递增的同胚，把 \([a, b]\) 映到 R 的区间 \([o, i]\) 上。（用超限归纳法证明：存在一个严格递增的同胚，把 \(A \cap [a, b]\) （把 A 等同于 X 中点 \((t, o)\) 所成之集）映到 \([o, i]\) 的一个子空间上。）
\end{enumerate}

¶ 13) 设 \(a\) 是一个无理数 \(>0\) 。对每个有理数 \(x\) ，令 \(f_{a}(x)\) 为区间 \([0, a]\) 中使 \(x - f_{a}(x)\) 是 \(a\) 的整数倍的那个实数。试证：\(f_{a}\) 是 \(\mathbf{Q}\) 到 \([0, a]\) 中的一个连续单射映射。试借助习题 9 推出：存在 \(\mathbf{Q}\) 到其自身的连续双射映射，其逆映射在任一点都不连续。

¶ 14) 设 X 是一个不只由一点组成的连通豪斯多夫拓扑空间，用 \(\Delta\) 表示 \(X \times X\) 的对角线。

\begin{enumerate}
\def\labelenumi{\alph{enumi})}
\item
试证：不存在把 \(\mathbb{C}\Delta\) 分成两个开集 \(\mathbf{D}_1, \mathbf{D}_2\) 的分划，使得 \(\mathbf{D}_1^{-1} = \mathbf{D}_1\) 且 \(\mathbf{D}_2^{-1} = \mathbf{D}_2\) 。{[}首先注意 \(\overline{\mathbf{D}_1(x)}\) 与 \(\overline{\mathbf{D}_2(x)}\) 对每个 \(x \in \mathbf{X}\) 都是连通的（利用第一章 § 11 习题 4 a)）；由此推出，若 \(y \in \mathbf{D}_1(x)\) ，则有 \(\overline{\mathbf{D}_2(x)} \times \{y\} \subset \mathbf{D}_1\) ，并证明这就蕴涵：对每个 \(x \in \mathbf{X}\) ，两集 \(\mathbf{D}_1(x)\) ， \(\mathbf{D}_2(x)\) 中必有一个为空；从而两集 \(\mathbf{D}_1, \mathbf{D}_2\) 中必有一个为空，这与假设矛盾。{]}由此推出：或者 \(\mathbb{C}\Delta\) 是连通的，或者 \(\mathbb{C}\Delta\) 恰有两个分支 A, B，使得 \(\mathbf{B} = \overline{\mathbf{A}}\) 。
\item
称 X 上的一个线性序与 X 的拓扑 \(\mathfrak{C}\) 相容，如果 \(\mathfrak{C}_0(\mathbf{X})\) 比 \(\mathfrak{C}\) 粗。试由 \(a\) ) 推出：X 上存在这样的线性序当且仅当 \(\mathbb{C}\Delta\) 不连通，并且这时恰有两个这样的线性序。{[}用 \(a\) ) 的记号，证明序关系必须是 \((x, y) \in \mathbf{A}\) 或 \((x, y) \in \mathbf{B} = \overline{\mathbf{A}}\) 。{]}试证：对这些序而言，各区间在拓扑 \(\mathfrak{C}\) 下都是连通的，并且它们是 X 关于 \(\mathfrak{C}\) 的仅有的连通子集。
\item
试证：若 X 在拓扑 \(\mathfrak{C}\) 下是局部连通或局部紧的，且 X 上存在一个与 \(\mathfrak{C}\) 相容的线性序，则必有 \(\mathfrak{C} = \mathfrak{C}_0(X)\) 。{[}若 X 局部连通，利用 b)。若 X 局部紧，通过考察邻域的边界，证明点 \(x \in X\) 的每个紧邻域也是 \(x\) 关于 \(\mathfrak{C}_0(X)\) 的邻域{]}。
\end{enumerate}

* d) 设 X 是 \(\mathbf{R}^2\) 的子空间，由 (o, o) 与所有点对 \((x, y)\) （满足 \(x > 0\) 且 \(y = \sin (\mathrm{i} / x)\) ）组成。试证：X 上存在一个线性序，它与 \(\mathbf{R}^2\) 的拓扑在 X 上诱导的拓扑相容，但该拓扑严格地细于 \(\mathcal{C}_0(X)\) 。 *

试给出一个连通豪斯多夫空间 X 的例子，它的任何一点都不具有连通邻域的基本系，而其上存在一个与 X 的拓扑相容的线性序{[}参看第一章，§ 11，习题 2 b){]}。

¶ 15) a) 设 X 是一个连通豪斯多夫空间，使得对每个 \(x \in X\) ，\(\complement\{x\}\) 恰有两个分支。试证：若 K 是这两个分支中的任何一个，则 \(K = K \cup \{x\}\) {[}第一章，§ 11，习题 4 a){]}。设 \(x, y\) 是 X 的两个不同的点，设 A 与 B 是 \(\complement\{x\}\) 的分支，并设 \(A', B'\) 是 \(\complement\{y\}\) 的分支。试证：两集 A, B 之一包含于 \(A'\) 或 \(B'\) 之中。

\begin{enumerate}
\def\labelenumi{\alph{enumi})}
\setcounter{enumi}{1}
\tightlist
\item
设 \(x_0\) 是 \(X\) 的一个点，并设 \(A(x_0), B(x_0)\) 是 \(\complement\{x_0\}\) 的分支。对每个 \(x \neq x_0\) ，\(\complement\{x\}\) 的两个分支中恰有一个或者包含于 \(A(x_0)\) 中，或者包含 \(A(x_0)\) ；把这个分支记作 \(A(x)\) 。试证：关系 \(A(x) \subset A(y)\) 是一个与 \(X\) 的拓扑相容的线性序关系（习题 14）。c) 把 b) 的结论推广到下述情形：\(\complement\{x\}\) 对每个 \(x \in X\) （除一个点 \(a \in X\) 外，对该点 \(\complement\{a\}\) 是连通的）都有两个分支。
\end{enumerate}

* \(d\) ) 设 \(\mathbf{X}\) 是 \(\mathbf{R}^2\) 的子空间，由点 \(a = (0, 1)\) ， \(b = (0, -1)\) 与所有点对 \((x, y)\) （满足 \(x \neq 0\) 且 \(y = \sin(1/x)\) ）组成。试证：\(\mathbf{X}\) 在拓扑 \(\mathfrak{G}\) （在 \(\mathbf{X}\) 上由 \(\mathbf{R}^2\) 的拓扑诱导）下是连通的，\(\mathbb{G}\{a\}\) 与 \(\mathbb{G}\{b\}\) 都是连通的，并且对每个点 \(x\) （ \(\mathbf{X}\) 中异于 \(a\) 与 \(b\) 的），\(\mathbb{G}\{x\}\) 恰有两个分支；但 \(\mathbf{X}\) 上不存在与拓扑 \(\mathfrak{G}_{*}\) 相容的线性序。

¶ 16) 设 X 是一个在两个不同点 \(a\) 与 \(b\) 之间完全不可约的连通豪斯多夫空间，也就是说，X 的连通子集中除 X 自身外没有一个同时含有 \(a\) 与 \(b\) 。试证：\(\mathbf{X}' = \mathbb{C}\{a, b\}\) 是连通的，并且对每个 \(x \in \mathbf{X}'\) ，\(x\) 在 X 中的余集恰有两个分支 \(\mathbf{A}(x)\) 与 \(\mathbf{B}(x)\) ，使得 \(a \in \mathbf{A}(x)\) ， \(b \in \mathbf{B}(x)\) ， \(a \notin \mathbf{B}(x)\) ， \(b \notin \mathbf{A}(x)\) 。{[}试利用第一章 § 11 习题 4 a) 证明 \(\mathbb{C}\{x\}\) 不能分成三个非空开集。{]}由此推出：X 上存在一个与 X 的拓扑相容的线性序（利用习题 15）。

¶ 17) 设 X 是一个连通豪斯多夫空间，使得每当 X 被三个异于 X 的非空连通集覆盖时，其中必有两个集的并异于 X。

\begin{enumerate}
\def\labelenumi{\alph{enumi})}
\item
试证：对每个点 \(x \in \mathbf{X}\) ，\(\mathbb{G}\{x\}\) 至多有两个分支（与习题 16 相同的方法）。
\item
试证：X 中至多有两个点，其余集是连通的。此外，若有两个不同的点 a, b 具有此性质，试证：对每个异于 a 或 b 的 x，a 与 b 属于 \(C\{x\}\) 的不同分支。
\item
试由 a) 与 b) 推出：X 上存在一个与 X 的拓扑相容的线性序（利用习题 15）。
\end{enumerate}

¶ 18) a) 设 X 是一个紧、连通、局部连通的空间，它在它的两个点 \(a\) , \(b\) 之间不可约（第二章，§ 4，习题 19）。设 \(x\) 是 X 中异于 \(a\) 或 \(b\) 的点。试证：\(\mathbb{C}\{x\}\) 恰有两个分支，且 \(a\) 与 \(b\) 属于 \(\mathbb{C}\{x\}\) 的不同分支。{[}按习题 16 的方式论证，首先证明 \(\mathbb{C}\{x\}\) 的分支不能多于两个。其次，若 V 是 \(x\) 的任一既不含 \(a\) 也不含 \(b\) 的连通闭邻域，令 \(\mathrm{A}_{\mathrm{V}}\) , \(\mathrm{B}_{\mathrm{V}}\) 分别是 \(a\) , \(b\) 在 \(\mathbb{C}\mathrm{V}\) 中的分支；利用第二章 § 4 习题 17 证明 \(\mathrm{A}_{\mathrm{V}} \neq \mathrm{B}_{\mathrm{V}}\) ；最后考察诸集 \(\mathrm{A}_{\mathrm{V}}\) （相应地 \(\mathrm{B}_{\mathrm{V}}\) ）当 V 取遍 \(x\) 在 X 中的所有连通闭邻域时的并。{]}由此推出：X 上存在一个线性序，使得 \(\mathfrak{C}_{0}(\mathbf{X})\) 就是 X 上给定的拓扑{[}习题 15 与 14 c){]}。

\begin{enumerate}
\def\labelenumi{\alph{enumi})}
\setcounter{enumi}{1}
\tightlist
\item
设 X 是一个紧连通空间，它在它的两个点 a, b 之间不可约。试证：若 X 的素成分（第二章，§ 4，习题 20）多于一个，则 X 的素成分所成的空间 \(X'\) （同处）在 a 与 b 的素成分之间不可约，并且 \(X'\) 上存在一个线性序，使得 \(X'\) 的拓扑等同于 \(\mathcal{G}_{0}(\mathbf{X}')\) 。 * 试给出一个 \(X'\) 异于 X 的例子{[}参看习题 15 d{]}。 *
\end{enumerate}

¶ 19) a) 设 X 是一个连通豪斯多夫空间，\(a, b\) 是 X 的两个不同的点，使得 X 的连通闭子集中除 X 自身外没有一个同时含有 \(a\) 与 \(b\) 。若 \(x\) 是 X 中异于 \(a\) 或 \(b\) 的点，则 \(\mathbb{C}\{x\}\) 不连通当且仅当：对每个 \(y \in \mathbb{C}\{x\}\) ，在 \(\mathbb{C}\{y\}\) 中存在一个在 X 中闭且含有 \(x\) 与两点 \(a, b\) 之一的连通子集。{[}为证条件必要，利用第一章 § 11 习题 4 a)；为证条件充分，令 A（相应地 B）是所有这样的 \(y \in X\) 所成之集：在 \(\mathbb{C}\{y\}\) 中存在一个在 X 中闭且含有 \(a\) 与 \(x\) （相应地 \(b\) 与 \(x\) ）的连通子集；证明 A 与 B 都非空、都是开集且互不相交，并且 \(\mathbb{C}\{x\} = A \cup B\) 。

\begin{enumerate}
\def\labelenumi{\alph{enumi})}
\setcounter{enumi}{1}
\tightlist
\item
设 X 是一个紧连通空间，它在它的两个点 \(a, b\) 之间不可约（第二章，§ 4，习题 19）。试证：若对 X 的每对不同的点 \(x, y\) ，X 中存在一个含有 \(y\) 与 \(a, b\) 之一但不含有 \(x\) 的紧连通集，则 X 上存在一个线性序，使得 \(\mathfrak{C}_{0}(\mathbf{X})\) 就是 X 上给定的拓扑（习题 15 与 16）。
\end{enumerate}

\begin{enumerate}
\def\labelenumi{\arabic{enumi})}
\setcounter{enumi}{19}
\item
在第一章 § 9 习题 23 的构造中，取 \(\mathbf{X}_0\) 为 R 的区间 {[}0, 1{]} 。试证：X 到其自身的同胚与 \(\mathbf{X}_0\) 到其自身的同胚相同，尽管 X 的拓扑严格地细于 \(\mathbf{X}_0\) 的拓扑{[}利用第一章 § 8 习题 20 b){]}。
\item
试证：若 \(r > 2\) ，则不存在 \(r\) 重传递的、\(\mathbf{R}\) 到其自身的同胚群（考察使 \(\mathbf{R}\) 的两点固定不动的子群）。
\end{enumerate}

